% The two topology types MPI knows about: a regular Cartesian grid, and a
% general graph where the edges are whatever neighbour relation you declare.
\documentclass[dvisvgm,border=3pt]{standalone}
\input{preamble.tex}
% Shared building block for the virtual-topology figures (slides 30, 31, 34,
% 35, 51): a 4x4 Cartesian grid of ranks with periodic wrap-around.
% Not a figure itself -- see tools/build_figs.py.

\tikzset{
  cdotp/.style={circle,fill=uibkorange!35,minimum size=4.6mm,inner sep=0pt},  % idle rank
  cdoth/.style={circle,fill=uibkorange,minimum size=4.6mm,inner sep=0pt},     % involved
  cdotc/.style={circle,fill=uibknavy,draw=blu,line width=0.5pt,
                minimum size=4.6mm,inner sep=0pt},                            % "me"
  clink/.style={draw=blu,line width=0.7pt},
  cwrap/.style={draw=blu,line width=0.7pt,dash pattern=on 2.5pt off 2.5pt},
}

\newcommand{\cartd}{0.78}   % spacing between ranks

% ---------------------------------------------------------------------------
% \carttopo{overrides}
%   Draws the grid, its links and the dashed periodic wrap-arounds, then the
%   ranks. Every rank is drawn pale; "overrides" is a comma list of
%   "col/row/style" (style: p pale, h highlighted, c centre) drawn on top.
%   Nodes are named t<col>-<row> so a figure can attach arrows to them.
% ---------------------------------------------------------------------------
\newcommand{\carttopo}[1]{%
  % nearest-neighbour links
  \foreach \c in {0,...,3}{\foreach \r in {0,...,2}{%
    \draw[clink] ({\c*\cartd},{-\r*\cartd}) -- ({\c*\cartd},{-\r*\cartd-\cartd});}}
  \foreach \c in {0,...,2}{\foreach \r in {0,...,3}{%
    \draw[clink] ({\c*\cartd},{-\r*\cartd}) -- ({\c*\cartd+\cartd},{-\r*\cartd});}}
  % periodic boundaries: each row and each column closes on itself
  \foreach \r in {0,...,3}{%
    \draw[cwrap,rounded corners=5pt]
      ({3*\cartd},{-\r*\cartd}) -- ({3*\cartd+0.42},{-\r*\cartd})
      -- ({3*\cartd+0.42},{-\r*\cartd+0.34}) -- ({-0.42},{-\r*\cartd+0.34})
      -- ({-0.42},{-\r*\cartd}) -- (0,{-\r*\cartd});}
  \foreach \c in {0,...,3}{%
    \draw[cwrap,rounded corners=5pt]
      ({\c*\cartd},{-3*\cartd}) -- ({\c*\cartd},{-3*\cartd-0.42})
      -- ({\c*\cartd+0.34},{-3*\cartd-0.42}) -- ({\c*\cartd+0.34},{0.42})
      -- ({\c*\cartd},{0.42}) -- ({\c*\cartd},0);}
  % ranks
  \foreach \c in {0,...,3}{\foreach \r in {0,...,3}{%
    \node[cdotp] (t\c-\r) at ({\c*\cartd},{-\r*\cartd}) {};}}
  \foreach \c/\r/\s in {#1}{%
    \node[cdot\s] (t\c-\r) at ({\c*\cartd},{-\r*\cartd}) {};}%
}

\begin{document}
\begin{tikzpicture}[x=1cm,y=1cm]

  \begin{scope}
    \carttopo{0/0/h, 1/0/h, 2/0/h, 3/0/h, 0/1/h, 1/1/h, 2/1/h, 3/1/h,
              0/2/h, 1/2/h, 2/2/h, 3/2/h, 0/3/h, 1/3/h, 2/3/h, 3/3/h}
  \end{scope}

  % an irregular graph: same idea, arbitrary edges
  \begin{scope}[xshift=5.6cm,yshift=-1.2cm]
    \foreach \n/\x/\y in {a/0.00/0.35, b/0.75/0.95, c/1.55/1.10, d/2.35/0.80,
                          e/3.10/1.05, f/3.80/0.55, g/0.55/-0.45, h/1.35/0.15,
                          i/2.05/-0.20, j/2.80/0.30, k/3.45/-0.30,
                          l/1.05/-1.05, m/1.95/-1.15, n/2.70/-0.85}{
      \node[cdoth] (\n) at (\x,\y) {};}
    \foreach \u/\v in {a/b, b/c, c/d, d/e, e/f, a/g, g/h, h/i, i/j, j/k, k/f,
                       b/h, c/h, c/i, d/i, d/j, e/j, g/l, l/m, m/n, n/k,
                       h/l, i/m, j/n, i/n, h/m}{
      \draw[clink] (\u) -- (\v);}
    % redraw the vertices so the edges sit underneath them
    \foreach \n/\x/\y in {a/0.00/0.35, b/0.75/0.95, c/1.55/1.10, d/2.35/0.80,
                          e/3.10/1.05, f/3.80/0.55, g/0.55/-0.45, h/1.35/0.15,
                          i/2.05/-0.20, j/2.80/0.30, k/3.45/-0.30,
                          l/1.05/-1.05, m/1.95/-1.15, n/2.70/-0.85}{
      \node[cdoth] at (\x,\y) {};}
  \end{scope}

\end{tikzpicture}
\end{document}
